\documentclass[11pt,letterpaper]{article}

\usepackage[utf8]{inputenc}
\usepackage[margin=0.9in]{geometry}
\usepackage{amsmath}
\usepackage{graphicx}
\usepackage{hyperref}
\usepackage{booktabs}
\usepackage{enumitem}
\usepackage{xcolor}
\usepackage{tcolorbox}
\usepackage{setspace}

% Compact spacing
\setlength{\parskip}{0.5em}
\setlist{noitemsep,topsep=0pt}

\title{\vspace{-2em}\textbf{CS 145 Project Proposals: Network Dynamics of AI Coding Tool Adoption}\vspace{-1em}}
\author{Team: [Your Name], [Partner 1], [Partner 2]}
\date{March 2026}

\begin{document}
\maketitle
\thispagestyle{empty}

%==============================================================================
% PROPOSAL 1
%==============================================================================
\section*{Proposal 1: Cascade Effects in AI Coding Tool Adoption}

\subsection*{The Idea}

We propose modeling AI coding tool adoption (Claude Code, GitHub Copilot, OpenAI Codex) as an \textbf{information cascade} through GitHub's developer collaboration network. When developers observe their collaborators successfully using AI coding assistants, they may be more likely to adopt these tools themselves---creating cascade dynamics similar to viral phenomena in social networks.

Our core research question: \textit{Does AI tool adoption spread through developer networks like an epidemic, and can we predict which developers will adopt based on their network position?}

We construct a \textbf{developer collaboration graph} $G = (V, E)$ where nodes are developers and edges connect those who contribute to shared repositories. Using timestamped adoption data from 5+ million AI-assisted commits across 600,000 repositories, we will:

\begin{enumerate}
    \item Reconstruct temporal adoption sequences for each AI tool
    \item Build cascade trees showing how adoption propagated through the network
    \item Fit Independent Cascade and Linear Threshold models to estimate transmission parameters
    \item Predict future adoptions based on current network state
\end{enumerate}

\subsection*{Motivation}

Understanding cascade dynamics in technology adoption has both theoretical and practical significance:

\begin{itemize}
    \item \textbf{Network Science}: Extends cascade models to a novel domain---professional developer networks---where ``infection'' represents tool adoption rather than information spread
    \item \textbf{Industry Applications}: AI companies can identify influential developers to seed adoption; organizations can predict which teams will adopt AI tools
    \item \textbf{Research Contribution}: First large-scale empirical analysis of how AI coding assistants spread through GitHub's collaboration network
\end{itemize}

The network perspective explains \textit{why} certain developers adopt faster than others, going beyond aggregate adoption curves.

\subsection*{End Product}

\begin{enumerate}
    \item \textbf{Interactive Visualization}: Cascade trees showing adoption spread for each AI tool
    \item \textbf{Calibrated Model}: Independent Cascade model with empirically estimated transmission rates ($\beta$) for each tool; comparison of cascade sizes and depths across tools
    \item \textbf{Predictive Analysis}: Identify developers in ``high-risk'' network positions likely to adopt within 6 months
    \item \textbf{Research Paper}: Empirical analysis comparing cascade dynamics across Claude, Copilot, Codex, and Jules
\end{enumerate}

\vspace{1em}
\hrule
\vspace{1em}

%==============================================================================
% PROPOSAL 2
%==============================================================================
\section*{Proposal 2: Disentangling Influence from Homophily in AI Tool Selection}

\subsection*{The Idea}

When two collaborating developers both use Claude Code, is it because one \textit{influenced} the other to adopt, or because similar developers independently make similar choices (\textit{homophily})? This is the fundamental identification problem in network analysis.

We propose applying rigorous causal inference techniques to our AI tool adoption dataset to decompose observed clustering into:
\begin{itemize}
    \item \textbf{Social Influence}: Developer $j$ adopts because neighbor $i$ convinced them
    \item \textbf{Homophily}: Developers with similar characteristics (language, domain) independently choose similar tools
    \item \textbf{Confounding}: Shared environment (same company, same open source project) drives both
\end{itemize}

Our methodology combines three identification strategies:

\begin{enumerate}
    \item \textbf{Temporal Analysis}: If influence is causal, neighbor adoption at $t-1$ predicts own adoption at $t$, controlling for covariates (programming language, repository topics, account age)
    \item \textbf{Matched Pair Design}: Compare adoption rates between developers exposed to adopting neighbors vs. matched controls without such exposure
    \item \textbf{Network Randomization}: Permutation tests that preserve network structure and marginal adoption rates to assess whether observed clustering exceeds chance
\end{enumerate}

\subsection*{Motivation}

This distinction has profound implications for technology diffusion strategy:
\begin{itemize}
    \item If \textbf{influence dominates}: Targeting network hubs can accelerate adoption exponentially
    \item If \textbf{homophily dominates}: Marketing should focus on developer segments rather than network position
    \item Understanding the mechanism enables better prediction of future technology waves
\end{itemize}

This project applies cutting-edge network causal inference methods (Aral et al. 2009, Shalizi \& Thomas 2011) to a timely domain with unique data.

\subsection*{End Product}

\begin{enumerate}
    \item \textbf{Decomposition}: Quantify the relative contributions of influence vs. homophily for each AI tool (e.g., ``60\% of Claude adoption clustering is attributable to peer influence'')
    \item \textbf{Tool Comparison}: Which tool shows strongest network effects? (Hypothesis: Claude Code, due to visible co-author attribution in commits)
    \item \textbf{Visualization}: Network diagrams with nodes colored by adoption source (influenced vs. independent)
    \item \textbf{Strategic Recommendations}: Evidence-based guidance for AI tool adoption strategies
\end{enumerate}

\vspace{1em}
\hrule
\vspace{0.5em}

%==============================================================================
% DATA SUMMARY
%==============================================================================
\section*{Data Available}

\begin{table}[h]
\centering
\small
\begin{tabular}{lrrl}
\toprule
\textbf{Tool} & \textbf{Events} & \textbf{Repos} & \textbf{Detection} \\
\midrule
Claude Code & 281K commits & 439K & Co-author trailer \\
GitHub Copilot & 1.08M PRs & 247K & \texttt{copilot/} branch \\
OpenAI Codex & 3.09M PRs & 249K & \texttt{codex/} branch \\
Google Jules & 622K commits & 70K & Bot author \\
\midrule
\textbf{Total} & \textbf{5.08M} & \textbf{$\sim$600K} & --- \\
\bottomrule
\end{tabular}
\end{table}

Data includes: developer usernames, timestamps, repository metadata (language, topics, stars), and owner location for geographic analysis. Collected Feb 2025--Jan 2026.

\end{document}
